% invariants.tex — Complete invariant set for the fee concentration insurance CFMM
% Kontrol: prove_cfmm_*, prove_fci_*, prove_ins_*

\section{Invariants}\label{sec:invariants}

This section contains all invariants for the fee concentration insurance system.
The invariant set is organized into five groups:
\begin{enumerate}
  \item CFMM-01 through CFMM-19: Base CFMM invariants (trading set, reserves, fees, ticks).
  \item CFMM-20 through CFMM-25: Replication invariants (payoff, reserves, canonical form).
  \item CFMM-26 through CFMM-28: Perpetual invariants (funding, jump safety).
  \item FCI-01 through FCI-13: Fee concentration index invariants (state, price mapping, perpetual funding).
  \item INS-01 through INS-10: Insurance CFMM invariants (borrow, settlement, utilization).
\end{enumerate}

% ══════════════════════════════════════════════════════════════════════════════
\subsection{Base CFMM Invariants (CFMM-01 through CFMM-19)}\label{sec:base-cfmm}
% ══════════════════════════════════════════════════════════════════════════════

\begin{invariant}{CFMM-01}{Trading set convexity}
% Kontrol: prove_cfmm_tradingSetConvexity
\[
  S = \{(u, y) : \psi(u, y) \geq L \cdot C_1\} \text{ is convex}
\]
The trading set is convex. Since $\psi(u,y) = y - (\pl^2/4)u$ is linear (hence concave),
the superlevel set is convex (Geometry paper \cite{angeris2023geometry}, Theorem~1).
\end{invariant}

\begin{invariant}{CFMM-02}{Reserve non-negativity}
% Kontrol: prove_cfmm_reserveNonNeg
\[
  u \geq 0 \quad\text{and}\quad y \geq 0
\]
Both virtual concentration reserve $u$ and USDC reserve $y$ are non-negative at all times.
\end{invariant}

\begin{invariant}{CFMM-03}{Invariant preservation on swap}
% Kontrol: prove_cfmm_invariantSwap
\[
  \{\psi(u, y) = L \cdot C_1\}
  \;\xrightarrow{\text{swap}(\Delta u)}\;
  \{\psi(u + \Delta u, y - \Delta y) = L \cdot C_1\}
\]
Every swap preserves the trading function invariant. From $\Delta y = (\pl^2/4)\Delta u$.
\end{invariant}

\begin{invariant}{CFMM-04}{Invariant preservation on mint}
% Kontrol: prove_cfmm_invariantMint
\[
  \{\psi(u, y) = L \cdot C_1\}
  \;\xrightarrow{\text{mint}(\Delta L)}\;
  \{\psi(u + \Delta u, y + \Delta y) = (L + \Delta L) \cdot C_1\}
\]
Minting new liquidity preserves the invariant at the new liquidity level
(\cref{prop:mint-invariant}).
\end{invariant}

\begin{invariant}{CFMM-05}{Invariant preservation on burn}
% Kontrol: prove_cfmm_invariantBurn
\[
  \{\psi(u, y) = L \cdot C_1\}
  \;\xrightarrow{\text{burn}(\Delta L)}\;
  \{\psi(u - \Delta u, y - \Delta y) = (L - \Delta L) \cdot C_1\}
\]
Burning liquidity preserves the invariant at the reduced liquidity level.
\end{invariant}

\begin{invariant}{CFMM-06}{Spot price consistency}
% Kontrol: prove_cfmm_spotPriceConsistency
\[
  p = \frac{\pl^2}{2}\sqrt{\frac{u}{\Lact}}
  = -\frac{\partial\psi/\partial u}{\partial\psi/\partial y} \cdot \frac{du}{dx}
\]
The spot price derived from reserves matches the marginal rate of substitution
(\cref{def:spot-price}).
\end{invariant}

\begin{invariant}{CFMM-07}{Price within tick range}
% Kontrol: prove_cfmm_priceInRange
\[
  \pl \leq p \leq \pu \quad\text{when } \Lact > 0
\]
The spot price stays within the active tick range $[\pl, \pu]$
when there is active liquidity.
\end{invariant}

\begin{invariant}{CFMM-08}{Liquidity additivity at tick crossing}
% Kontrol: prove_cfmm_tickCrossing
\[
  \Lact' = \Lact + \Lnet
\]
At each tick crossing, active liquidity updates by adding the signed net liquidity
of the crossed tick. This is O(1) because $\psi$ is 1-homogeneous in $L$
(\cref{thm:linearity}).
\end{invariant}

\begin{invariant}{CFMM-09}{Liquidity non-negativity}
% Kontrol: prove_cfmm_liquidityNonNeg
\[
  L \geq 0 \quad\text{and}\quad \Lact \geq 0
\]
Total and active liquidity are non-negative.
\end{invariant}

\begin{invariant}{CFMM-10}{No-arbitrage (monotone reserves)}
% Kontrol: prove_cfmm_noArbitrage
\[
  p' > p \;\longrightarrow\; u(p') > u(p) \;\land\; y(p') > y(p)
\]
Both reserves are strictly increasing in price (\cref{prop:reserve-monotone}).
Under Framing~C, rising concentration absorbs more virtual exposure and
accumulates more USDC.
\end{invariant}

\begin{invariant}{CFMM-11}{Fee accumulator monotonicity}
% Kontrol: prove_cfmm_feeAccMonotone
\[
  g^{\text{borrow}}_{n+1} \geq g^{\text{borrow}}_n
\]
The borrow fee growth accumulator is non-decreasing (borrow rate $\geq \premfactor > 0$).
\end{invariant}

\begin{invariant}{CFMM-12}{Fee accumulator per-position}
% Kontrol: prove_cfmm_feeAccPerPosition
\[
  \text{fees}_i = L_i \cdot (g^{\text{borrow}} - g^{\text{borrow}}_{\text{baseline},i})
\]
Per-position borrow fees use the Synthetix difference pattern \cite{batog2018scalable}.
\end{invariant}

\begin{invariant}{CFMM-13}{LP value bounded}
% Kontrol: prove_cfmm_lpValueBounded
\[
  0 \leq V(p) = \frac{\pu^2 - p^2}{\pl^2} \leq \frac{\pu^2 - \pl^2}{\pl^2}
\]
LP value is bounded between 0 (at $p = \pu$, full loss) and the initial
collateral (at $p = \pl$, no loss).
\end{invariant}

\begin{invariant}{CFMM-14}{Swap output positive}
% Kontrol: prove_cfmm_swapOutputPositive
\[
  \Delta x > 0 \;\longrightarrow\; \Delta y > 0
  \quad\text{and}\quad
  \Delta x < 0 \;\longrightarrow\; \Delta y < 0
\]
Swaps in either direction produce correctly-signed outputs.
\end{invariant}

\begin{invariant}{CFMM-15}{Slippage non-negative}
% Kontrol: prove_cfmm_slippageNonNeg
\[
  \frac{\Delta y}{\Delta x} = \frac{\pl^2}{4}(2x + \Delta x) \geq \frac{\pl^2}{4} \cdot 2x = p
\]
The effective price is always at least the spot price (natural slippage from
the $u = x^2$ transform). Equality only for infinitesimal trades.
\end{invariant}

\begin{invariant}{CFMM-16}{Coordinate transform bijection}
% Kontrol: prove_cfmm_coordTransformBijection
\[
  u = x^2 \;\longleftrightarrow\; x = \sqrt{u} \quad\text{for } x \geq 0
\]
The $u \leftrightarrow x$ mapping is a bijection on the non-negative reals.
One \texttt{FixedPointMathLib.sqrt} call per swap.
\end{invariant}

\begin{invariant}{CFMM-17}{1-homogeneity}
% Kontrol: prove_cfmm_oneHomogeneity
\[
  \psi(\lambda u, \lambda y) = \lambda \psi(u, y) \quad\forall\, \lambda > 0
\]
The trading function is 1-homogeneous, enabling O(1) tick-based liquidity
aggregation (\cref{thm:linearity}).
\end{invariant}

\begin{invariant}{CFMM-18}{Concavity}
% Kontrol: prove_cfmm_concavity
\[
  \nabla^2 \psi =
  \begin{pmatrix} 0 & 0 \\ 0 & 0 \end{pmatrix}
\]
$\psi(u,y) = y - (\pl^2/4)u$ is linear, hence (weakly) concave.
The Hessian is zero. The trading set is convex.
\end{invariant}

\begin{invariant}{CFMM-19}{Single-sided deposit consistency}
% Kontrol: prove_cfmm_singleSidedDeposit
\[
  \Delta L = \frac{\Delta y \cdot \pl^2}{\pu^2 + p^2}
  \quad\text{and}\quad
  \Delta u = \Delta L \cdot \frac{4p^2}{\pl^4}
\]
Single-sided USDC deposits correctly compute $\Delta L$ and virtual $\Delta u$
(\cref{def:single-sided}). Only USDC enters the contract.
\end{invariant}

% ══════════════════════════════════════════════════════════════════════════════
\subsection{Replication Invariants (CFMM-20 through CFMM-25)}\label{sec:replication}
% ══════════════════════════════════════════════════════════════════════════════

Source: Angeris et al.\ \cite{angeris2021replicating,angeris2023geometry}.

\begin{invariant}{CFMM-20}{Replication accuracy}
% Kontrol: prove_cfmm_replicationAccuracy
\[
  \{V(p)\}
  \;\xrightarrow{\text{evolve}(p \to p')}\;
  \{|\text{LP\_value}(p') - V(p')| \leq \varepsilon\}
\]
The CFMM LP position replicates the target payoff $V(p) = (\pu^2 - p^2)/\pl^2$
within error $\varepsilon$ (determined by discretization and fee leakage).
\end{invariant}

\begin{invariant}{CFMM-21}{LP payoff concavity}
% Kontrol: prove_cfmm_payoffConcavity
\[
  V''(p) = -\frac{2}{\pl^2} < 0 \quad\forall\, p \in [\pl, \pu]
\]
The LP payoff is strictly concave on the active domain.
Required for Angeris replication theorem \cite{angeris2021replicating}.
\end{invariant}

\begin{invariant}{CFMM-22}{Reserve-payoff consistency}
% Kontrol: prove_cfmm_reservePayoffConsistency
\[
  x(p) = -V'(p) = \frac{2p}{\pl^2}
  \quad\text{and}\quad
  y(p) = V(p) + p \cdot x(p) = \frac{\pu^2 + p^2}{\pl^2}
\]
Reserves are the derivatives and Legendre transform of the payoff
(\cref{def:reserves}).
\end{invariant}

\begin{invariant}{CFMM-23}{Boundary reserves}
% Kontrol: prove_cfmm_boundaryReserves
\begin{align*}
  p = \pl &:\quad x = 2/\pl,\; y = (\pu^2 + \pl^2)/\pl^2,\; V = (\pu^2 - \pl^2)/\pl^2 \\
  p = \pu &:\quad x = 2\pu/\pl^2,\; y = 2\pu^2/\pl^2,\; V = 0
\end{align*}
At the lower boundary, full collateral remains. At the upper boundary, collateral
is fully drained (max loss for underwriter, max payout for PLP).
\end{invariant}

\begin{invariant}{CFMM-24}{Price-reserve consistency}
% Kontrol: prove_cfmm_priceReserveConsistency
\[
  p = \frac{\pl^2}{2} \cdot \frac{x}{L}
  \quad\Longleftrightarrow\quad
  x(p) = \frac{2pL}{\pl^2}
\]
The price derived from the trading function marginal rate matches the
reserve curve parameterization (\cref{def:spot-price}).
\end{invariant}

\begin{invariant}{CFMM-25}{Canonical form}
% Kontrol: prove_cfmm_canonicalForm
\[
  \psi(u, y) = y - \frac{\pl^2}{4}\, u
\]
is nondecreasing in $y$, nonincreasing in $u$, concave, and 1-homogeneous.
This is the unique canonical trading function for this CFMM
(Geometry paper \cite{angeris2023geometry}, Theorem~1).
\end{invariant}

% ══════════════════════════════════════════════════════════════════════════════
\subsection{Perpetual Invariants (CFMM-26 through CFMM-28)}\label{sec:perpetual-invariants}
% ══════════════════════════════════════════════════════════════════════════════

Source: Angeris et al.\ \cite{angeris2022perpetuals}.

\begin{invariant}{CFMM-26}{Funding rate convergence}
% Kontrol: prove_cfmm_perpetual_fundingConvergence
\[
  \{|\pmark - \pindex| > \delta\}
  \;\xrightarrow{\text{funding}(\Delta t)}\;
  \{|\pmark' - \pindex'| < |\pmark - \pindex|\}
\]
The bidirectional funding mechanism drives mark-index convergence
(\cref{prop:convergence}). Sensitivity $\alpha > 0$ controls speed.
\end{invariant}

\begin{invariant}{CFMM-27}{Funding rate bounds}
% Kontrol: prove_cfmm_perpetual_fundingBounds
\[
  |\rfunding_n| \leq \alpha + \lambda \cdot |\Delta\pindex|_{\max}
\]
The funding rate (including jump premium) is bounded. The bound depends on
the sensitivity $\alpha$, jump intensity $\lambda$, and maximum observable
$\pindex$ change between callbacks.
\end{invariant}

\begin{invariant}{CFMM-28}{Jump-adjusted value bound}
% Kontrol: prove_cfmm_perpetual_jumpSafety
\[
  \{\DeltaPlus \text{ jumps from } \delta \text{ to } \delta'\}
  \;\longrightarrow\;
  \{|\Delta\text{value}| \leq |V(\delta') - V(\delta)|\}
\]
Discrete jumps in $\DeltaPlus$ at liquidity events produce value changes
bounded by the payoff difference (\cref{thm:jump-safety}).
\end{invariant}

% ══════════════════════════════════════════════════════════════════════════════
\subsection{Fee Concentration State Invariants (FCI-01 through FCI-08)}\label{sec:fci-state}
% ══════════════════════════════════════════════════════════════════════════════

\begin{invariant}{FCI-01}{Index boundedness}
% Kontrol: prove_fci_indexBoundedness
\[
  0 \leq \AT \leq 1
\]
The fee concentration index is bounded in $[0,1]$ at all times.
$\AT = 0$ when no positions have been removed. $\AT = 1$ is the theoretical
maximum (single JIT monopoly).
\end{invariant}

\begin{invariant}{FCI-02}{Theta sum non-negativity}
% Kontrol: prove_fci_thetaSumNonNeg
\[
  \ThetaSum \geq 0
\]
The aggregate turnover rate $\ThetaSum = \sum_k \thetapos_k$ is non-negative
since each $\thetapos_k = 1/\blocklife_k > 0$.
\end{invariant}

\begin{invariant}{FCI-03}{Position count non-negativity}
% Kontrol: prove_fci_posCountNonNeg
\[
  \PosCount \geq 0
\]
The active position count is a non-negative integer.
\end{invariant}

\begin{invariant}{FCI-04}{Null lower bound}
% Kontrol: prove_fci_nullLowerBound
\[
  \{\PosCount > 0\} \;\longrightarrow\; \{\AT \geq \ATnull\}
\]
The real fee concentration index is always at least as large as the Ma--Crapis
equal-share null (\cref{prop:null-lower-bound}). Equivalently, $\DeltaPlus \geq 0$.
\end{invariant}

\begin{invariant}{FCI-05}{Deviation non-negativity}
% Kontrol: prove_fci_deviationNonNeg
\[
  \DeltaPlus = \max(0,\, \AT - \ATnull) \geq 0
\]
By construction ($\max$ with zero) and by \cref{prop:null-lower-bound}.
\end{invariant}

\begin{invariant}{FCI-06}{Deviation upper bound}
% Kontrol: prove_fci_deviationUpperBound
\[
  \DeltaPlus < 1
\]
Since $\AT \leq 1$ and $\ATnull \geq 0$, the deviation satisfies $\DeltaPlus \leq 1$.
Strict inequality holds when $\PosCount \geq 1$ (at least one position implies $\ATnull > 0$).
This ensures the price mapping $p = \DeltaPlus/(1-\DeltaPlus)$ is well-defined.
\end{invariant}

\begin{invariant}{FCI-07}{Co-primary consistency}
% Kontrol: prove_fci_coPrimaryConsistency
\[
  \ATnull = \sqrt{\frac{\ThetaSum}{\PosCount^2}}
  \quad\text{and}\quad
  \DeltaPlus = \max(0,\, \AT - \ATnull)
\]
The derived quantities $\ATnull$ and $\DeltaPlus$ are deterministic functions of
the three stored state variables $(\AT, \ThetaSum, \PosCount)$.
No additional state is needed.
\end{invariant}

\begin{invariant}{FCI-08}{Liquidity event atomicity}
% Kontrol: prove_fci_liquidityEventAtomicity
\[
  \{\texttt{afterRemoveLiquidity}\}
  \;\longrightarrow\;
  \{\AT \text{ updated} \;\land\; \ThetaSum \text{ updated} \;\land\; \PosCount \text{ updated}\}
\]
All three co-primary state variables update atomically at each liquidity event.
No intermediate state is observable where some variables are updated and others are not.
\end{invariant}

% ══════════════════════════════════════════════════════════════════════════════
\subsection{Price Mapping Invariants (FCI-09 through FCI-11)}
% ══════════════════════════════════════════════════════════════════════════════

\begin{invariant}{FCI-09}{Price non-negativity}
% Kontrol: prove_fci_priceNonNeg
\[
  p = \frac{\DeltaPlus}{1 - \DeltaPlus} \geq 0
\]
Since $\DeltaPlus \geq 0$ (FCI-05) and $1 - \DeltaPlus > 0$ (FCI-06).
\end{invariant}

\begin{invariant}{FCI-10}{Price monotonicity}
% Kontrol: prove_fci_priceMonotonicity
\[
  \DeltaPlus_1 < \DeltaPlus_2
  \;\longrightarrow\;
  p(\DeltaPlus_1) < p(\DeltaPlus_2)
\]
The price mapping is strictly increasing (\cref{prop:price-monotone}).
Higher concentration deviation $\to$ higher insurance price.
\end{invariant}

\begin{invariant}{FCI-11}{Price-deviation invertibility}
% Kontrol: prove_fci_priceInvertibility
\[
  \DeltaPlus = \frac{p}{1 + p}
  \quad\text{is the unique inverse of}\quad
  p = \frac{\DeltaPlus}{1 - \DeltaPlus}
\]
The mapping is a bijection from $[0,1) \to [0,\infty)$.
\end{invariant}

% ══════════════════════════════════════════════════════════════════════════════
\subsection{FCI Perpetual Funding Invariants (FCI-12 through FCI-13)}
% ══════════════════════════════════════════════════════════════════════════════

Source: Angeris et al.\ \cite{angeris2022perpetuals}.

\begin{invariant}{FCI-12}{Funding rate convergence (FCI)}
% Kontrol: prove_fci_fundingConvergence
\[
  \{|\pmark - \pindex| > \delta\}
  \;\longrightarrow\; \text{funding}(\text{liquidityEvent})
  \;\longrightarrow\; \{|\pmark' - \pindex'| < |\pmark - \pindex|\}
\]
The funding mechanism drives mark-index convergence (\cref{prop:convergence}).
\end{invariant}

\begin{invariant}{FCI-13}{Jump-adjusted value bound (FCI)}
% Kontrol: prove_fci_jumpBound
\[
  \{\DeltaPlus \text{ jumps from } \delta \text{ to } \delta'\}
  \;\longrightarrow\;
  \{|\Delta\text{value}| \leq |V(\delta') - V(\delta)|\}
\]
Discrete jumps in $\DeltaPlus$ at liquidity events do not cause value changes
exceeding the payoff difference (\cref{thm:jump-safety}).
\end{invariant}

% ══════════════════════════════════════════════════════════════════════════════
\subsection{Insurance CFMM Invariants (INS-01 through INS-10)}\label{sec:insurance-cfmm}
% ══════════════════════════════════════════════════════════════════════════════

\begin{invariant}{INS-01}{Collateral solvency}
% Kontrol: prove_ins_collateralSolvency
\[
  y \geq L \cdot \frac{\pu^2 + p^2}{\pl^2} \quad\text{(USDC reserve covers all LP positions)}
\]
The USDC reserve is sufficient to cover the numeraire reserve requirement
at the current price. Underwriters cannot lose more than their deposit
(\cref{def:collateral-value}).
\end{invariant}

\begin{invariant}{INS-02}{Utilization bound}
% Kontrol: prove_ins_utilizationBound
\[
  0 \leq \Ua = \frac{L_{\text{borrowed}}}{L_{\text{total}}} \leq 1
\]
Utilization is bounded in $[0,1]$. $L_{\text{borrowed}} \leq L_{\text{total}}$
is enforced by reverting borrows that would exceed available liquidity.
\end{invariant}

\begin{invariant}{INS-03}{Borrow rate positivity}
% Kontrol: prove_ins_borrowRatePositivity
\[
  \rborrow(\Ua) \geq \premfactor > 0 \quad\forall\, \Ua \in [0,1]
\]
The borrow rate is always at least $\premfactor$ (the base premium floor).
Underwriters always earn a positive rate (\cref{prop:borrow-properties}).
\end{invariant}

\begin{invariant}{INS-04}{Borrow rate monotonicity}
% Kontrol: prove_ins_borrowRateMonotonicity
\[
  \Ua_1 < \Ua_2 \;\longrightarrow\; \rborrow(\Ua_1) \leq \rborrow(\Ua_2)
\]
The borrow rate is non-decreasing in utilization. Higher utilization
means higher cost, incentivizing new underwriter deposits or PLP exit
(\cref{prop:borrow-properties}).
\end{invariant}

\begin{invariant}{INS-05}{Separate accumulator accounting}
% Kontrol: prove_ins_separateAccumulators
\[
  g^{\text{borrow}} \geq 0
  \quad\text{(always positive)}
  \qquad
  g^{\text{funding}} \in \mathbb{R}
  \quad\text{(bidirectional)}
\]
The borrow and funding accumulators are maintained separately
(\cref{def:borrow-accumulator,def:funding-accumulator}).
The borrow accumulator is non-negative (PLP always pays underwriter).
The funding accumulator can be positive or negative (bidirectional).
Conflating them breaks funding convergence \cite{angeris2022perpetuals}.
\end{invariant}

\begin{invariant}{INS-06}{Max-price tracking}
% Kontrol: prove_ins_maxPriceTracking
\[
  p_{\max,i} = \max_{t \in [\text{mint}_i, \text{now}]} \pindex(t)
  \quad\text{and}\quad
  p_{\max,i}' \geq p_{\max,i}
\]
Each PLP position stores the running maximum of $\pindex$ since mint.
Updated at every hook callback when $\pindex > p_{\max,i}$.
This value is monotone non-decreasing.
\end{invariant}

\begin{invariant}{INS-07}{Payout formula}
% Kontrol: prove_ins_payoutFormula
\[
  \text{payout}_i = \text{premium}_i \times \left[\left(\frac{p_{\max,i}}{\pl}\right)^2 - 1\right]^{+}
\]
The PLP payout at exit follows the power-2 payoff. The plus operator
$[\cdot]^+$ ensures non-negative payouts. The premium is defined as
$\text{premium}_i = \premfactor \cdot \text{lifetime\_fees}_i
+ \text{accrued\_borrow}_i + \text{accrued\_funding}_i$.
\end{invariant}

\begin{invariant}{INS-08}{Payout cap}
% Kontrol: prove_ins_payoutCap
\[
  \text{payout}_i \leq \text{premium}_i \times \left[\left(\frac{\pu}{\pl}\right)^2 - 1\right]
\]
The payout is capped by the tick range. Since $p_{\max,i} \leq \pu$ within
the CLAMM range, the payout cannot exceed the cap. No global cap parameter
is needed---the per-position tick range enforces bounded loss for underwriters.
\end{invariant}

\begin{invariant}{INS-09}{Initialization invariant}
% Kontrol: prove_ins_initInvariant
\[
  \psi(u_0, y_0) = L_0 \cdot \frac{\pu^2}{\pl^2}
\]
At initialization (first deposit), the invariant holds with
$C_1 = \pu^2/\pl^2$ (\cref{prop:init-invariant}).
\end{invariant}

\begin{invariant}{INS-10}{Exit settlement conservation}
% Kontrol: prove_ins_exitSettlement
\[
  \{\texttt{exit}_i\}
  \;\longrightarrow\;
  \{L_{\text{borrowed}}' = L_{\text{borrowed}} - L_i\}
  \;\land\;
  \{\text{net}_i = \max(\text{payout}_i - \text{premium}_i, 0)\}
\]
At PLP exit: borrowed liquidity returns to the pool (utilization decreases),
and the PLP receives the net settlement (payout minus premium, floored at zero).
If payout $<$ premium, the excess stays in reserves as underwriter profit.
\end{invariant}
